% id: 23-33
% book: 베이직쎈 확률과 통계 · 02 중복조합과 이항정리
% section: 개념 06 방정식의 해의 개수
% passage: 다음 방정식의 양의 정수인 해의 개수를 구하시오.
% figure: none
% answer: 
% answer_source: 
% variant_level: 0 (원본 전사)
% 이 파일은 build.mjs 가 items.json 에서 만든다. 고칠 때는 items.json 을 고친다.
\smash{\raisebox{1.5mm}{\dmkichul{베이직쎈 확률과 통계 23쪽 33번}}}
$x+y+z=6$
\dmhint{$x$, $y$, $z$가 양의 정수이므로 $x-1=a$, $y-1=b$, $z-1=c$로 놓으면 $a$, $b$, $c$는 음이 아닌 정수이다. $x+y+z=6$에서\\ $a+b+c=\blank$ \cond{$a$, $b$, $c$는 음이 아닌 정수}\\ 따라서 구하는 해의 개수는 $a$, $b$, $c$의 3개에서 \blank개를 택하는 중복조합의 수와 같으므로\\ ${}_{\text{\scriptsize\setlength{\fboxsep}{1.5pt}\framebox[1.5em]{\rule{0pt}{1.6ex}}}}\mathrm{H}_{\text{\scriptsize\setlength{\fboxsep}{1.5pt}\framebox[1.5em]{\rule{0pt}{1.6ex}}}}={}_{\text{\scriptsize\setlength{\fboxsep}{1.5pt}\framebox[1.5em]{\rule{0pt}{1.6ex}}}}\mathrm{C}_3=\blank$}
